\begin{table}[h]
    \centering
    \begin{tabular}{|c|p{13cm}|c|}
        \hline
        1 & To‘g‘ri chiziqli tekis harakat & $\checkmark$ \\ \hline
        2 & To‘g‘ri chiziqli tekis harakatni grafik ravishda tasvirlash   & $\checkmark$\\ \hline
        3 & Tekis o`zgaruvchan harakat. Tezlanish & $\checkmark$ \\ \hline
        4 & Tekis o‘zgaruvchan harakatda ko‘chish   & $\checkmark$\\ \hline
        5 &   Tekis o‘zgaruvchan harakatni grafik ravishda tasvirlash & $\checkmark$\\ \hline
        6 &  Harakatning nisbiyligi  & $\checkmark$\\ \hline
        7 & Notekis harakat   & $\checkmark$\\ \hline
        8 & Kuch. Nyutonning birinchi qonuni   & $\checkmark$\\ \hline
        9 & Massa va zichlik   &$\checkmark$ \\ \hline
        10 & Butun olam tortishish qonuni   & $\checkmark$\\ \hline
        11 & Og‘irlik kuchi   &$\checkmark$ \\ \hline
        12 &  Gorizontal otish  & $\checkmark$\\ \hline
        13 & Burchak ostida otish   &$\checkmark$ \\ \hline
        14 & Ishqalanish koeffitsiyenti   & $\checkmark$\\ \hline
        15 & Ishqalanishli qiya tekislik   &$\checkmark$ \\ \hline
        16 & Bir nechta kuch ta’siridagi harakat   & $\checkmark$\\ \hline
        17 &  Jism impulsi. Impulsning saqlanish qonuni  & $\checkmark$\\ \hline
        18 & Impulsning o‘zgarishi va o‘rtacha kuch   &$\checkmark$ \\ \hline
        19 &  Ish va energiya  &$\checkmark$ \\ \hline
        20 &  Kinetik energiya. Ish va kinetik energiyaning o‘zgarishi  &$\checkmark$ \\ \hline
        21 & Potensial energiya. Ish va potensial energiyaning o‘zgarishi   &$\checkmark$ \\ \hline
        22 &  Mexanik energiyaning saqlanish qonuni  &$\checkmark$ \\ \hline
        23 & Mexanik energiyaning ichki energiyaga o‘tishi   &$\checkmark$ \\ \hline
        24 &  Mexanik energiyaning tashqi kuchlar ta’sirida o`zgarishi  &$\checkmark$ \\ \hline
        25 &  O‘zgarmas quvvat. FIK  & $\checkmark$\\ \hline
        26 &  Bosim. Paskal qonuni  &$\checkmark$ \\ \hline
        27 & Gidrostatik bosim   &$\checkmark$ \\ \hline
        28 &  Tutash idishlar  &$\checkmark$ \\ \hline
        29 & Atmosfera bosimi   & $\checkmark$\\ \hline
        30 &  Gidravlik press  &$\checkmark$ \\ \hline
        31 &  Prujinali va matematik mayatnik  &$\checkmark$ \\ \hline
        32 &  Prujinali mayatnik  &$\checkmark$ \\ \hline
        33 & Matematik mayatnik   &$\checkmark$ \\ \hline
        34 &  Mexanik to‘lqinlar  & $\checkmark$\\ \hline
        35 &   MKN & $\checkmark$\\ \hline
        36 &  Ichki energiya  & $\checkmark$ \\ \hline
        37 &  Elektr maydon. Elektr maydon kuchlanganligi  &$\checkmark$ \\ \hline
        38 &  Potensiallar farqi  &$\checkmark$ \\ \hline
        39 & Ekvipotensial sirtlar   &$\checkmark$ \\ \hline
        40 & Zanjirning bir qismi uchun Om qonuni. Qarshilik.   & $\checkmark$\\ \hline
        41 &  Ampermetr va voltmetrga qo`shimcha qarshilik ulash  &$\checkmark$ \\ \hline
        42 &  O`zgarmas tokning ishi va quvvati  & $\checkmark$\\ \hline
        43 & Joul Lens qonuni   &$\checkmark$ \\ \hline
        44 & Tok manbalarini ketma-ket va parallel ulash   &$\checkmark$ \\ \hline
        45 & Elektromagnit induksiya qonuni.  Lens qoidasi.   &$\checkmark$ \\ \hline
        46 & Tebranishlar konturida zaryad, kuchlanish va tok kuchi   &$\checkmark$ \\ \hline
        47 &  Ideal tebranishlar konturi uchun Tomson formulasi  & $\checkmark$\\ \hline
        48 & Tebranishlar konturida energiya.   &$\checkmark$ \\ \hline
        49 & O‘zgaruvchan tok zanjirida aktiv  qarshilik, kondensator va induktiv g‘altak   & $\checkmark$\\ \hline
        50 & Yorug‘likning qaytishi.   &$\checkmark$ \\ \hline
        51 & Yorug‘likning sinishi   &$\checkmark$ \\ \hline
        52 & Yorug‘likning to‘lqin tabiati   &$\checkmark$ \\ \hline
        53 & To‘la ichki qaytish   &$\checkmark$ \\ \hline
        54 & Linzalar   & $\checkmark$\\ \hline
        55 & Linzalar beradigan tasvirlar   &$\checkmark$ \\ \hline
        56 &  Linza formulasi  & $\checkmark$\\ \hline
        57 &  Optik asboblar  &$\checkmark$ \\ \hline
        3 &    &$\checkmark$ \\ \hline
        3 &    &$\checkmark$ \\ \hline
        3 &    &$\checkmark$ \\ \hline
        3 &    & $\checkmark$\\ \hline
        3 &    &$\checkmark$ \\ \hline
        3 &    &$\checkmark$ \\ \hline
        3 &    &$\checkmark$ \\ \hline
        3 &    &$\checkmark$ \\ \hline
        3 &    & $\checkmark$\\ \hline
        3 &    &$\checkmark$ \\ \hline
        3 &    & $\checkmark$\\ \hline
        3 &    &$\checkmark$ \\ \hline
        3 &    &$\checkmark$ \\ \hline
        3 &    &$\checkmark$ \\ \hline
        3 &    &$\checkmark$ \\ \hline
        3 &    & $\checkmark$\\ \hline
        3 &    &$\checkmark$ \\ \hline


        
    \end{tabular}
    
    \label{tab:placeholder}
\end{table}